% Pursuing Stacks (cotes 134-2 to 134-8): archive cote 134-3, site page 43 (PDF page 44), typescript page 43.
% typescript. Draft compared: ps2.tex:2116-2192.
% AI-transcribed by Claude Opus 5.5 (claude-opus-5-5, adaptive thinking, effort high), Message Batches API, batch ps1 (msgbatch_01Rv5hCaNTtUt4JHx1yqoHPL), 2026-11-12; not yet checked by eye.
\scanpage{134-3/43}
\mspage{43}
33.   For the time being,  I'll use the word ``test category'' \unsure{with} the mea-%
ning of last week,  and refer to categories satisfying the conditions
T 1 to  T 3  as \uline{strict}  test categories. (NB A is supposed to be essential\inserted[typed]{ly small in any case.)}

First of all,  the conditions T 1) and T 3) are \uline{necessary} for A to be
a test category.  For  T 1)  this is just the ``concrete'' criterion (iii)
of yesterday (page 38),  when B is the final element of (Cat).  For T \over{\illegible{1}}{3}),
we get even a \uline{canonical} \uline{choice} for $I,e_o,e_1$ ,  namely starting with the
\xout{x} ``universal'' choice in (Cat) :
\[ \underline{I} = \hand{\Delta}_1 \ , \quad \underline{e}_o \text{ and } \underline{e}_1 \text{ the two subobjects of } \hand{\Delta}_1 \text{ in} \]
(Cat) (or in (Ord)) corresponding to the
two unique sections of $\underline{I}$ over the final
element $\underline{e}$ of (Cat)\hand{,}
i.e. in terms of $\underline{I}$ as an ordered set $\hand{\{}0\hand{\}} \hand{\longrightarrow} \hand{\{}1\hand{\}}$ ,  $\underline{e}_o$ and $\underline{e}_1$ are just the
two subobjects defined by the two vertices $\hand{\{}0\hand{\}}$ and $\hand{\{}1\hand{\}}$ (viewed as defining
two one-point ordered subsets of $\underline{I} = \hand{\Delta}_1$).  We now apply $j_A$ to get
\[ I = j_A(\underline{I}) = j_A(\hand{\Delta}_1) \ , \quad e_i = j_A(\underline{e}_i) \quad \text{for } i = 0,1 \ . \]
As  $\underline{I} \hand{\to} \underline{e}$  is a weak equivalence so is $I \hand{\to} e_{A^{\hand{\wedge}}} = j_A(\underline{e})$ ,  and hence (as $e_{A^{\hand{\wedge}}}$
is \unsure{asphéric|aspheric} by T 1) I is aspheric.  As $\underline{e}_o \hand{\wedge} \underline{e}_1 = \emptyset_{(\mathrm{Cat})}$  and $j_A$ commutes
with inverse limits,  and with sums (and in particular  \unsure{transforms} the
\margin{($A^{\wedge}$) circled}
initial element \xout{of} $\emptyset$ of (Cat) in the initial element $\emptyset$ of $A^{\hand{\wedge}}$ ), it follows
$e_o \hand{\wedge} e_1 = \emptyset$ .

It's worth while having a look at what this \deleted{\illegible{1}}\inserted{object} just constructed is
like.  For this end,  let's note first that for any category\inserted{C}\hand{,} we have
a canonical bijection,  fonctorial in C
\[ \mathrm{Hom}(C, \hand{\Delta}_1) \hand{\xrightarrow{\sim}} \mathrm{Crib}(C) \hand{\simeq} \text{set of all subobjects of } e_{C^{\wedge}} \ , \]
by associating to any functor $f\colon\ C \hand{\to} \hand{\Delta}_1$  the full subcategory $f^{-1}(\underline{e}_o)$ of
\margin{(sieve) circled}
C,  which clearly is a ``crible''\inserted{[=sieve]} in C.  Thus we get,  for a $\hand{\in}$ Ob A
\[ I(a) \hand{\simeq} \mathrm{Crib}(A_{/a}) \quad \hand{(}\text{subobjects of a in } A^{\wedge}\hand{)}. \]
M\over{\illegible{1}}{o}re generally,  we deduce from this,  for any F $\hand{\in}$ Ob $A^{\wedge}$:
\[ \mathrm{Hom}(F,I) = \hand{\Gamma}(I_F = I\mathrm{x}F \ / \ F) = \text{ set of all subobjects of } F\hand{,} \]
in other words, the object I is defined intrinsically in the topos $A^{\wedge}$\hand{/}(up
to unique isomorphism)  as the ``\uline{Lawvere element}''  representing the functor
$F \hand{\mapsto}$ set of subobjects of F ,  generalizing the functor  $F \hand{\mapsto \mathfrak{P}}(F)$ in the
category of sets (namely sheaves over the one-point topos)\hand{,} \inserted{(represented by the two-point set $\{0,1\}$.)} The  condition
$e_o \hand{\wedge} e_1 = \emptyset_{\hand{\mathcal{A}}}$  is automatic in any topos $\hand{\mathcal{A}}$, provided $\hand{\mathcal{A}}$ is not the ``empty
topos''  namely corresponding to a category of sheaves equivalent to a one-%
point category. \deleted{Thus The condition  T 3' )  deduced from  T 3) by being
specific about the present choice of $(I,e_o,e_1)$ ,  is equivalent to T 3)
when T 1) and T 2) are satisfied,  and \xout{has the advantage} means simply that}
